% GENERATED FILE. Edit canonical lessons, then run npm run generate:books.
% canonical-source-hash: fnv1a64:d2d24981876e3d66
% canonical-lessons: UR-C210-nai, UR-C210-dhobi, UR-C210-mochi, UR-C210-lohar, UR-C210-barhai

\chapter{Barber, Washerman, Cobbler, Blacksmith, Carpenter}
\label{ch:ur-barber-washerman-cobbler-blacksmith-carpenter}
\hlchaptermodality{210}

\begin{chapteropening}
\textbf{By the end of this chapter:} \emph{I can say a barber, a washerman, a cobbler, a blacksmith and a carpenter.}

Complete the last lesson of chapter 210: I can say a barber, a washerman, a cobbler, a blacksmith and a carpenter.
\end{chapteropening}

\section[\texorpdfstring{\ur{نائی} (nāī)}{نائی (nāī)}]{\ur{نائی} (nāī) — a barber}

\label{lesson:UR-C210-nai}



\begin{quote}
Before the new one: say the Urdu for a bill, then the Urdu for an expense.
\end{quote}

\subsection*{You'll want to know: \ur{نائی}}
\textbf{\ur{نائی}} — \emph{nāī} — \textquotedblleft{}a barber\textquotedblright{}.

\begin{cousinweb}
\textbf{\ur{نائی}} (\emph{nāī}) is a tadbhava, inherited and reshaped by sound change. CDIAL lists it under Sanskrit \emph{nāpita}, \textquotedblleft{}barber\textquotedblright{}, and lists \emph{nāū}, \textquotedblleft{}barber\textquotedblright{}, beside it under \emph{snāpaka}, \textquotedblleft{}bath attendant\textquotedblright{} (CDIAL 13792).
\end{cousinweb}

\begin{grammarlens}[title={What you've built}]
One more word for this chapter.
\end{grammarlens}

\subsection*{Your turn}
\begin{itemize}
\raggedright
  \item \emph{Say it:} \emph{nāī}
  \item \emph{Say it:} \emph{nāī}, once more
  \item \emph{Say it:} say \emph{nāī}
  \item \emph{Recall:} say the Urdu for a bill, then the Urdu for an expense, then say \emph{nāī} again
\end{itemize}

\begin{tcolorbox}[breakable,colback=teal!4,colframe=teal!35!black,title={Before you move on}]
What does \ur{نائی} mean? (\textquotedblleft{}A barber\textquotedblright{}.) Say it once more.
\end{tcolorbox}
\section[\texorpdfstring{\ur{دھوبی} (dhobī)}{دھوبی (dhobī)}]{\ur{دھوبی} (dhobī) — a washerman}

\label{lesson:UR-C210-dhobi}



\begin{quote}
Before the new one: say the Urdu for an expense, then the Urdu for a barber.
\end{quote}

\subsection*{You'll want to know: \ur{دھوبی}}
\textbf{\ur{دھوبی}} — \emph{dhobī} — \textquotedblleft{}a washerman\textquotedblright{}.

\begin{cousinweb}
\textbf{\ur{دھوبی}} (\emph{dhobī}) is a tadbhava, a word inherited and reshaped by sound change. It comes from \emph{*dhauvati}, \textquotedblleft{}washes\textquotedblright{}, a form CDIAL reconstructs from its descendants (CDIAL 6886). CDIAL notes that the forms with \emph{-b-} spread from here into every language around.
\end{cousinweb}

\begin{grammarlens}[title={What you've built}]
One more word for this chapter.
\end{grammarlens}

\subsection*{Your turn}
\begin{itemize}
\raggedright
  \item \emph{Say it:} \emph{dhobī}
  \item \emph{Say it:} \emph{dhobī}, once more
  \item \emph{Say it:} say \emph{dhobī}
  \item \emph{Recall:} say the Urdu for an expense, then the Urdu for a barber, then say \emph{dhobī} again
\end{itemize}

\begin{tcolorbox}[breakable,colback=teal!4,colframe=teal!35!black,title={Before you move on}]
What does \ur{دھوبی} mean? (\textquotedblleft{}A washerman\textquotedblright{}.) Say it once more.
\end{tcolorbox}
\section[\texorpdfstring{\ur{موچی} (mochī)}{موچی (mochī)}]{\ur{موچی} (mochī) — a cobbler}

\label{lesson:UR-C210-mochi}



\begin{quote}
Before the new one: say the Urdu for a barber, then the Urdu for a washerman.
\end{quote}

\subsection*{You'll want to know: \ur{موچی}}
\textbf{\ur{موچی}} — \emph{mochī} — \textquotedblleft{}a cobbler\textquotedblright{}.

\begin{cousinweb}
\textbf{\ur{موچی}} (\emph{mochī}) comes from \emph{mocika}, \textquotedblleft{}tanner and cobbler\textquotedblright{}, built on \emph{*moca}, \textquotedblleft{}shoe\textquotedblright{}, which CDIAL takes from an Iranian word, Pahlavi \emph{mōč}, \textquotedblleft{}shoe\textquotedblright{} (CDIAL 10349).
\end{cousinweb}

\begin{grammarlens}[title={What you've built}]
One more word for this chapter.
\end{grammarlens}

\subsection*{Your turn}
\begin{itemize}
\raggedright
  \item \emph{Say it:} \emph{mochī}
  \item \emph{Say it:} \emph{mochī}, once more
  \item \emph{Say it:} say \emph{mochī}
  \item \emph{Recall:} say the Urdu for a barber, then the Urdu for a washerman, then say \emph{mochī} again
\end{itemize}

\begin{tcolorbox}[breakable,colback=teal!4,colframe=teal!35!black,title={Before you move on}]
What does \ur{موچی} mean? (\textquotedblleft{}A cobbler\textquotedblright{}.) Say it once more.
\end{tcolorbox}
\section[\texorpdfstring{\ur{لوہار} (lohār)}{لوہار (lohār)}]{\ur{لوہار} (lohār) — a blacksmith}

\label{lesson:UR-C210-lohar}



\begin{quote}
Before the new one: say the Urdu for a washerman, then the Urdu for a cobbler.
\end{quote}

\subsection*{You'll want to know: \ur{لوہار}}
\textbf{\ur{لوہار}} — \emph{lohār} — \textquotedblleft{}a blacksmith\textquotedblright{}.

\begin{cousinweb}
\textbf{\ur{لوہار}} (\emph{lohār}) is a tadbhava, a word inherited from Sanskrit and reshaped by sound change. It comes from \emph{lohakāra}, \textquotedblleft{}iron-worker\textquotedblright{} (CDIAL 11159). It is built from \emph{loha}, \textquotedblleft{}iron\textquotedblright{}, and \emph{kāra}, \textquotedblleft{}maker\textquotedblright{}.
\end{cousinweb}

\begin{grammarlens}[title={What you've built}]
One more word for this chapter.
\end{grammarlens}

\subsection*{Your turn}
\begin{itemize}
\raggedright
  \item \emph{Say it:} \emph{lohār}
  \item \emph{Say it:} \emph{lohār}, once more
  \item \emph{Say it:} say \emph{lohār}
  \item \emph{Recall:} say the Urdu for a washerman, then the Urdu for a cobbler, then say \emph{lohār} again
\end{itemize}

\begin{tcolorbox}[breakable,colback=teal!4,colframe=teal!35!black,title={Before you move on}]
What does \ur{لوہار} mean? (\textquotedblleft{}A blacksmith\textquotedblright{}.) Say it once more.
\end{tcolorbox}
\section[\texorpdfstring{\ur{بڑھئی} (baṛhaī)}{بڑھئی (baṛhaī)}]{\ur{بڑھئی} (baṛhaī) — a carpenter}

\label{lesson:UR-C210-barhai}



\begin{quote}
Before the new one: say the Urdu for a cobbler, then the Urdu for a blacksmith.
\end{quote}

\subsection*{You'll want to know: \ur{بڑھئی}}
\textbf{\ur{بڑھئی}} — \emph{baṛhaī} — \textquotedblleft{}a carpenter\textquotedblright{}.

\begin{cousinweb}
\textbf{\ur{بڑھئی}} (\emph{baṛhaī}) is a tadbhava, a word inherited from Sanskrit and reshaped by sound change. It comes from \emph{vardhaki}, \textquotedblleft{}carpenter\textquotedblright{} (CDIAL 11375).
\end{cousinweb}

\begin{grammarlens}[title={What you've built}]
One more word for this chapter.
\end{grammarlens}

\subsection*{Your turn}
\begin{itemize}
\raggedright
  \item \emph{Say it:} \emph{baṛhaī}
  \item \emph{Say it:} \emph{baṛhaī}, once more
  \item \emph{Say it:} say \emph{baṛhaī}
  \item \emph{Recall:} say the Urdu for a cobbler, then the Urdu for a blacksmith, then say \emph{baṛhaī} again
\end{itemize}

\begin{tcolorbox}[breakable,colback=teal!4,colframe=teal!35!black,title={Before you move on}]
What does \ur{بڑھئی} mean? (\textquotedblleft{}A carpenter\textquotedblright{}.) Say it once more.
\end{tcolorbox}
